\documentclass[11pt,a4paper]{article}
\usepackage[a4paper,top=15mm,bottom=15mm,left=16mm,right=16mm]{geometry}
\usepackage[T1]{fontenc}
\usepackage{lmodern}
\usepackage{amsmath,amssymb}
\usepackage{enumitem}
\usepackage{multicol}
\usepackage{array}
\usepackage{xcolor}

\setlength{\parindent}{0pt}
\setlength{\parskip}{4pt}
\renewcommand{\arraystretch}{1.2}
\newcolumntype{C}[1]{>{\centering\arraybackslash}m{#1}}
\setlist[enumerate,1]{label=\textbf{Q\arabic*.}, leftmargin=*, itemsep=4pt, topsep=2pt}
\setlist[enumerate,2]{label=(\alph*), leftmargin=*, itemsep=3pt, topsep=2pt}

\begin{document}

\begin{center}
{\Large\bfseries Worked Solutions: From Times Tables to the $n$th Term}
\end{center}

\section*{Part 1: Times tables}

\textbf{Q1.} The first six terms are:
\begin{center}\small
\begin{tabular}{|c|*{6}{c|}}\hline
Table & \(\times1\) & \(\times2\) & \(\times3\) & \(\times4\) & \(\times5\) & \(\times6\)\\\hline
2 & 2 & 4 & 6 & 8 & 10 & 12\\
5 & 5 & 10 & 15 & 20 & 25 & 30\\
3 & 3 & 6 & 9 & 12 & 15 & 18\\
7 & 7 & 14 & 21 & 28 & 35 & 42\\
12 & 12 & 24 & 36 & 48 & 60 & 72\\
15 & 15 & 30 & 45 & 60 & 75 & 90\\\hline
\end{tabular}
\end{center}

\textbf{Q2.} Decimal times tables:
\begin{center}\small
\begin{tabular}{|c|*{6}{c|}}\hline
Table & \(\times1\) & \(\times2\) & \(\times3\) & \(\times4\) & \(\times5\) & \(\times6\)\\\hline
\(0.1\) & \(0.1\) & \(0.2\) & \(0.3\) & \(0.4\) & \(0.5\) & \(0.6\)\\
\(0.5\) & \(0.5\) & \(1\) & \(1.5\) & \(2\) & \(2.5\) & \(3\)\\
\(1.5\) & \(1.5\) & \(3\) & \(4.5\) & \(6\) & \(7.5\) & \(9\)\\
\(0.25\) & \(0.25\) & \(0.5\) & \(0.75\) & \(1\) & \(1.25\) & \(1.5\)\\\hline
\end{tabular}
\end{center}

\textbf{Q3.} Fraction times tables, in simplest form:
\begin{center}\small
\begin{tabular}{|c|*{6}{c|}}\hline
Table & \(\times1\) & \(\times2\) & \(\times3\) & \(\times4\) & \(\times5\) & \(\times6\)\\\hline
\(\frac12\) & \(\frac12\) & \(1\) & \(\frac32\) & \(2\) & \(\frac52\) & \(3\)\\
\(\frac14\) & \(\frac14\) & \(\frac12\) & \(\frac34\) & \(1\) & \(\frac54\) & \(\frac32\)\\
\(\frac13\) & \(\frac13\) & \(\frac23\) & \(1\) & \(\frac43\) & \(\frac53\) & \(2\)\\
\(\frac34\) & \(\frac34\) & \(\frac32\) & \(\frac94\) & \(3\) & \(\frac{15}{4}\) & \(\frac92\)\\\hline
\end{tabular}
\end{center}

\textbf{Q4.} Percentage times tables:
\begin{center}\small
\begin{tabular}{|c|*{6}{c|}}\hline
Table & \(\times1\) & \(\times2\) & \(\times3\) & \(\times4\) & \(\times5\) & \(\times6\)\\\hline
\(10\%\) & \(10\%\) & \(20\%\) & \(30\%\) & \(40\%\) & \(50\%\) & \(60\%\)\\
\(25\%\) & \(25\%\) & \(50\%\) & \(75\%\) & \(100\%\) & \(125\%\) & \(150\%\)\\
\(5\%\) & \(5\%\) & \(10\%\) & \(15\%\) & \(20\%\) & \(25\%\) & \(30\%\)\\
\(12.5\%\) & \(12.5\%\) & \(25\%\) & \(37.5\%\) & \(50\%\) & \(62.5\%\) & \(75\%\)\\\hline
\end{tabular}
\end{center}

\textbf{Q5.} Negative times tables:
\begin{center}\small
\begin{tabular}{|c|*{6}{c|}}\hline
Table & \(\times1\) & \(\times2\) & \(\times3\) & \(\times4\) & \(\times5\) & \(\times6\)\\\hline
\(-2\) & \(-2\) & \(-4\) & \(-6\) & \(-8\) & \(-10\) & \(-12\)\\
\(-5\) & \(-5\) & \(-10\) & \(-15\) & \(-20\) & \(-25\) & \(-30\)\\
\(-0.5\) & \(-0.5\) & \(-1\) & \(-1.5\) & \(-2\) & \(-2.5\) & \(-3\)\\
\(-\frac14\) & \(-\frac14\) & \(-\frac12\) & \(-\frac34\) & \(-1\) & \(-\frac54\) & \(-\frac32\)\\\hline
\end{tabular}
\end{center}

\textbf{Q6.}
\begin{enumerate}[label=(\alph*)]
\item \(9,18,\mathbf{27},36\). This is the 9 times table.
\item \(0.7,\mathbf{1.4},2.1,2.8\). This is the \(0.7\) times table.
\item \(\frac15,\mathbf{\frac25},\frac35,\mathbf{\frac45}\). This is the \(\frac15\) times table.
\item \(\mathbf{-6},-12,-18,\mathbf{-24}\). This is the \(-6\) times table.
\end{enumerate}

\textbf{Q7.} Complete table:
\begin{center}\small
\begin{tabular}{|p{5.8cm}|*{5}{C{1.1cm}|}C{1.8cm}|}\hline
Number of tickets, minutes or hours & 1 & 2 & 3 & 4 & 5 & Table\\\hline
Cost of cinema tickets at \pounds7.50 each
& \pounds7.50 & \pounds15.00 & \pounds22.50 & \pounds30.00 & \pounds37.50 & \(7.5n\)\\\hline
Height above sea level of a submarine going down \(30\)\,m each minute
& \(-30\) & \(-60\) & \(-90\) & \(-120\) & \(-150\) & \(-30n\)\\\hline
Battery used by a phone that uses \(15\%\) of its charge each hour
& \(15\%\) & \(30\%\) & \(45\%\) & \(60\%\) & \(75\%\) & \(15n\%\)\\\hline
\end{tabular}
\end{center}
The times tables are therefore \(7.5n\) (or the \pounds7.50 times table), \(-30n\), and \(15n\%\).

\section*{Part 2: Substituting into expressions}

\textbf{Q8.}
\begin{center}\small
\begin{tabular}{|c|*{6}{c|}c|}\hline
Expression & \(n=1\) & \(n=2\) & \(n=3\) & \(n=4\) & \(n=5\) & \(n=10\) & Times table\\\hline
\(2n\) & 2 & 4 & 6 & 8 & 10 & 20 & 2\\
\(3n\) & 3 & 6 & 9 & 12 & 15 & 30 & 3\\
\(7n\) & 7 & 14 & 21 & 28 & 35 & 70 & 7\\
\(12n\) & 12 & 24 & 36 & 48 & 60 & 120 & 12\\\hline
\end{tabular}
\end{center}

\textbf{Q9.}
\begin{center}\small
\begin{tabular}{|c|*{6}{c|}c|}\hline
Expression & \(n=1\) & \(n=2\) & \(n=3\) & \(n=4\) & \(n=5\) & \(n=10\) & Times table\\\hline
\(0.1n\) & \(0.1\) & \(0.2\) & \(0.3\) & \(0.4\) & \(0.5\) & \(1\) & \(0.1\)\\
\(0.5n\) & \(0.5\) & \(1\) & \(1.5\) & \(2\) & \(2.5\) & \(5\) & \(0.5\)\\
\(\frac14n\) & \(\frac14\) & \(\frac12\) & \(\frac34\) & \(1\) & \(\frac54\) & \(\frac52\) & \(\frac14\)\\
\(5n\%\) & \(5\%\) & \(10\%\) & \(15\%\) & \(20\%\) & \(25\%\) & \(50\%\) & \(5\%\)\\
\(-2n\) & \(-2\) & \(-4\) & \(-6\) & \(-8\) & \(-10\) & \(-20\) & \(-2\)\\
\(-0.5n\) & \(-0.5\) & \(-1\) & \(-1.5\) & \(-2\) & \(-2.5\) & \(-5\) & \(-0.5\)\\\hline
\end{tabular}
\end{center}

\textbf{Q10.}
\begin{enumerate}[label=(\alph*)]
\item The 100th term of \(7n\) is \(7\times100=700\).
\item The 100th term of \(0.5n\) is \(0.5\times100=50\).
\item The 20th term of \(-3n\) is \(-3\times20=-60\).
\item No. If \(7n=60\), then
\[
n=\frac{60}{7}=8\frac47,
\]
which is not a position number \(1,2,3,\ldots\). Equivalently, \(7\times8=56\) and \(7\times9=63\), so 60 lies between consecutive terms and is not a term of the sequence.
\end{enumerate}

\section*{Part 3: Matching sequences to their \(n\)th terms}

\textbf{Q11.}
\begin{enumerate}[label=(\alph*)]
\item Matching sequences to their \(n\)th terms:
\begin{center}\small
\begin{tabular}{|c|c|}\hline
Sequence & \(n\)th term\\\hline
\(6,12,18,24,\ldots\) & \(6n\)\\
\(0.4,0.8,1.2,1.6,\ldots\) & \(0.4n\)\\
\(-4,-8,-12,-16,\ldots\) & \(-4n\)\\
\(\frac14,\frac12,\frac34,1,\ldots\) & \(0.25n\) or \(\frac14n\)\\
\(20\%,40\%,60\%,80\%,\ldots\) & \(20n\%\)\\
\(-0.5,-1,-1.5,-2,\ldots\) & \(-0.5n\)\\\hline
\end{tabular}
\end{center}
The two \(n\)th terms left over are \(n+6\) and \(0.5n\).

\item First four terms:
\[
n+6:\quad 7,\ 8,\ 9,\ 10,\ldots
\]
\[
0.5n:\quad 0.5,\ 1,\ 1.5,\ 2,\ldots
\]
\end{enumerate}

\section*{Part 4: Picture sequences}

\textbf{Q12.}
\begin{center}\small
\begin{tabular}{|l|*{6}{c|}c|}\hline
Picture sequence & Pattern 1 & Pattern 2 & Pattern 3 & Pattern 4 & Pattern 5 & Pattern 10 & \(n\)th term\\\hline
(a) Dots & 3 & 6 & 9 & 12 & 15 & 30 & \(3n\)\\
(b) Separate squares of matchsticks & 4 & 8 & 12 & 16 & 20 & 40 & \(4n\)\\
(c) Separate hexagons of matchsticks & 6 & 12 & 18 & 24 & 30 & 60 & \(6n\)\\\hline
\end{tabular}
\end{center}
For (a), pattern \(n\) has \(n\) rows of 3 dots, giving \(3n\) dots. For (b), each square uses 4 matchsticks, so \(n\) squares use \(4n\). For (c), each hexagon uses 6 matchsticks, so \(n\) hexagons use \(6n\).

\textbf{Q13.}
\begin{center}\small
\begin{tabular}{|l|*{6}{c|}c|}\hline
Picture sequence & Pattern 1 & Pattern 2 & Pattern 3 & Pattern 4 & Pattern 5 & Pattern 10 & \(n\)th term\\\hline
(a) Zigzag matchsticks & 2 & 4 & 6 & 8 & 10 & 20 & \(2n\)\\
(b) Tooth triangles & 3 & 6 & 9 & 12 & 15 & 30 & \(3n\)\\
(c) Cross of dots & 4 & 8 & 12 & 16 & 20 & 40 & \(4n\)\\\hline
\end{tabular}
\end{center}
The zigzag has \(2n\) matchsticks. Each triangle uses 3 matchsticks and the triangles do not share sides, so \(n\) triangles use \(3n\) matchsticks. The cross has four arms of \(n\) dots, so it has \(4n\) dots.

\section*{Part 5: Matching and drawing picture sequences}

\textbf{Q14.} Matching picture sequences to their \(n\)th terms:
\begin{center}\small
\begin{tabular}{|c|c|c|}\hline
Picture sequence & Counting & \(n\)th term\\\hline
A: dice showing five & 5 dots per die, so \(5n\) dots & \(5n\)\\
B: tables with chairs & 4 chairs per table, so \(4n\) chairs & \(4n\)\\
C: star of dots & 3 arms of \(n\) dots, so \(3n\) dots & \(3n\)\\
D: flowers of dots & centre \(+\) 6 petals \(=7\) dots per flower, so \(7n\) & \(7n\)\\
E: rows of 2 dots & \(n\) rows of 2, so \(2n\) dots & \(2n\)\\
F: houses of matchsticks & 6 matchsticks per house, so \(6n\) & \(6n\)\\\hline
\end{tabular}
\end{center}
The three \(n\)th terms left over are
\[
5n+1,\qquad n+4,\qquad 3n-2.
\]

\textbf{Q15.} Sample answers. In each case, draw patterns 1, 2 and 3 with the stated numbers of objects.
\begin{enumerate}[label=(\alph*)]
\item \(3n\): use \(3,6,9\) objects in patterns 1, 2 and 3. For example, draw 3 dots, then 6 dots, then 9 dots in matching arrangements.
\item \(5n\): use \(5,10,15\) objects in patterns 1, 2 and 3.
\item \(8n\): use \(8,16,24\) objects in patterns 1, 2 and 3.
\item \(\frac12n\): the numbers are
\[
\frac12,\quad 1,\quad \frac32.
\]
This can be drawn if half-objects are allowed, for example half a matchstick, one matchstick, and one and a half matchsticks. With whole objects only, pattern 1 cannot be shown exactly.
\item One possible choice is \(4n\). Then draw \(4,8,12\) objects in patterns 1, 2 and 3.
\item No. The term \(-2n\) gives
\[
-2,\ -4,\ -6,\ldots
\]
for \(n=1,2,3,\ldots\), and a picture cannot contain a negative number of objects.
\end{enumerate}

\section*{Part 6: Always, sometimes or never?}

\textbf{Q16--Q27.}
\begin{center}\small
\begin{tabular}{|p{6.4cm}|p{2.2cm}|p{6.2cm}|}\hline
\textbf{Statement} & \textbf{Answer} & \textbf{Reason}\\\hline
Q16. A term of the 6 times table is also a term of the 3 times table.
& Always
& Every multiple of 6 is a multiple of 3, since \(6n=3(2n)\).\\\hline
Q17. A term of the 3 times table is also a term of the 6 times table.
& Sometimes
& True for multiples such as 6; false for 3, because 3 is not a multiple of 6.\\\hline
Q18. A term of the \(25\%\) times table is more than \(100\%\).
& Sometimes
& \(25n\%\) is more than \(100\%\) only when \(n>4\). For example, \(n=1\) gives \(25\%\), but \(n=5\) gives \(125\%\).\\\hline
Q19. Adding 2 to every term of a times table gives another times table.
& Never
& Adding 2 to \(an\) gives \(an+2\). For example, \(3,6,9,\ldots\) becomes \(5,8,11,\ldots\), which is not a standard times table.\\\hline
Q20. A term of the sequence \(2n+1\) is even.
& Never
& \(2n\) is even, so \(2n+1\) is always odd.\\\hline
Q21. \(-2n\) is less than \(-n\).
& Always
& For \(n\ge1\), \(2n>n\), so \(-2n<-n\).\\\hline
Q22. \(3n\) is greater than \(n+2\).
& Sometimes
& \(3n>n+2\) is equivalent to \(2n>2\), i.e. \(n>1\). It is true for \(n\ge2\), but false for \(n=1\) because \(3=3\).\\\hline
Q23. In the sequence \(3n+1\), doubling the position doubles the term.
& Never
& The term at position \(2n\) is \(3(2n)+1=6n+1\), while twice the term at position \(n\) is \(2(3n+1)=6n+2\). These differ by 1.\\\hline
Q24. A number that is a term of both \(3n\) and \(5n\) is also a term of \(15n\).
& Always
& Common multiples of 3 and 5 are exactly the multiples of 15.\\\hline
Q25. The number of chairs in a pattern of these tables is even.
& Always
& Each table has 4 chairs, so \(n\) tables have \(4n\) chairs, which is even.\\\hline
Q26. The number of dots on a row of these dice is even.
& Sometimes
& Each die shows 5 dots, so \(n\) dice show \(5n\) dots. This is even when \(n\) is even and odd when \(n\) is odd.\\\hline
Q27. A pattern of these hexagons uses exactly 100 matchsticks.
& Never
& Each hexagon uses 6 matchsticks, so a pattern of \(n\) hexagons uses \(6n\) matchsticks. But \(100\) is not divisible by 6.\\\hline
\end{tabular}
\end{center}

\textbf{Q28.} Sample statements:
\begin{itemize}
\item Always: Every term of the sequence \(6n\) is a multiple of 3.
\item Sometimes: A term of the sequence \(4n\) is greater than 10. This is true for \(n=3\) because \(4\times3=12\), but false for \(n=1\) because \(4\times1=4\).
\item Never: A term of the sequence \(5n\) is negative. Since \(n=1,2,3,\ldots\), every term \(5n\) is positive.
\end{itemize}

\end{document}